\documentclass[10pt,a4paper]{scrartcl}		% Ersatz für Klasse "article" 
\parskip 12pt
\parindent 0pt
\usepackage{makeidx}					% Stichwortverzeichnis
\makeindex
\usepackage[totoc]{idxlayout}
% \makeglossary
\usepackage[utf8]{inputenc}				% Umlaute in Quelldatei UTF8-kodiert
\usepackage[T1]{fontenc}
\usepackage{footnpag}					% Nummerierung auf jeder Seite neu beginnen
\usepackage[hyphens]{url}
\usepackage{listings}
\lstloadlanguages{[Visual]Basic}
\lstset{language=C, basicstyle=\ttfamily\small, commentstyle=\itshape }
\usepackage{numprint}
\usepackage{booktabs}
\usepackage{amssymb}					% $\checkmark$
\usepackage[colorlinks=true,
urlcolor=blue,
pdfpagelabels,
pdfstartview = FitH,
pdftitle={XD-2031 Benutzerhandbuch},
pdfauthor={Nils Eilers},
bookmarksopen = true,
bookmarksnumbered = true,
linkcolor = blue,
plainpages = false,
hypertexnames = false,
citecolor = black] {hyperref}
\usepackage{eurosym}
\usepackage{hyperref}
\usepackage{tikz}
\usetikzlibrary{trees}
\usepackage{verbatim}
\usepackage{graphicx}
\usepackage[ngerman]{babel}				% Umlaute "a "s etc.
\RequirePackage[ngerman=ngerman-x-latest]{hyphsubst}

\hyphenation{zwi-schen}


% **************   SERIF   **************
%\usepackage{mathptmx} % Times Roman
%usepackage{fouriernc} % New Century Schoolbook
\usepackage{tgschola} % TEX Gyre Schola (based on New Century Schoolbook)
%\usepackage[scaled]{DejaVuSerif}

% **************   SANS SERIF   **************
%\usepackage[scaled=.90] {helvet}
\usepackage{tgadventor} % TEX Gyre Adventor (based on URW Gothic, an Avantgarde clone)
%\usepackage{tgheros} % TEX Gyre Heros (based on URW Nimbus Sans)
%\usepackage{arev} % Arev is a version of Bitstream Vera Sans, disigned for slide presentations

% **************   MONO   **************
%\usepackage{luximono}
\usepackage{beramono}

\makeatletter
\def\url@foostyle{%
  \@ifundefined{selectfont}{\def\UrlFont{\sf}}{\def\UrlFont{\small\ttfamily}}}
\makeatother

\urlstyle{foo}

\newcommand{\beginsf}[1]{#1\begin{addmargin}[0.5cm]{0cm}}
\newcommand{\sfend}[1]{\end{addmargin}}

% -------------------------------------------------------------------------
% Global variables for LaTeX documents
% -------------------------------------------------------------------------

% Internet links
\newcommand{\fwbinaries}{\url{http://xd2031.petsd.net/firmware.php}}
\newcommand{\osxbinaries}{\url{http://xd2031.petsd.net/serverosx.php}}
\newcommand{\srcdownloads}{\url{https://github.com/fachat/XD2031/tags}}


\begin{document}

\title{Benutzerhandbuch}
\subtitle{XD-2031 Serial line filesystem server}
\author{}
\maketitle\thispagestyle{empty}

\clearpage

\tableofcontents
\clearpage

\section{Allgemeines}

\subsection{Verwendungszweck dieser Software}
XD-2031 ist Software, die einen modernen Mikrocontroller für einen
historischen Commodore-Computer wie ein Diskettenlaufwerk erscheinen lässt

Sie dient dem Datenaustausch zwischen alten und modernen Computern, der
ohne weiteres nicht mehr möglich ist. 
Moderne Computer können die Datenträger der alten nicht mehr lesen,
die alten aber auch nicht die der neuen und sie haben auch keine gemeinsame 
Schnittstelle, über die sie Daten austauschen könnten.

XD-2031 schließt diese Lücke mithilfe eines Adapters zwischen den
damals verwendeten Bus-Systemen und der aktuellen USB-Schnittstelle.

Damit kann man beispielsweise...
\begin{itemize}
\item alte Daten und Programme auf einem neuen PC sichern
\item aus dem Internet geladene Programme auf dem Commodore ausprobieren
\item mit modernen Werkzeugen Software für alte Rechner entwickeln
\end{itemize}

Dabei sind hier ausschließlich die 8-Bit Commodore-Computer gemeint, die 
über einen parallelen IEEE-488-Bus \index{IEEE-488-Bus} oder seriellen
	\index{IEC-Bus} \index{CBM-Bus|see{IEC-Bus}} IEC-Bus\footnote
	{Diese proprietäre serielle Variante des IEEE-488-Busses wird
	auch als CBM-Bus bezeichnet, da er tatsächlich \textit{nicht} durch die
	\textit{International Electrotechnical Commission} standardisiert
	wurde.}
verfügen:
\begin{itemize}
\item C64, C128, VC20, +4
\item CBM 3032, 
PET 2001\footnote{
	Nur mit BASIC 2, erkennbar an der Einschaltmeldung 
	\texttt{\#\#\# COMMODORE BASIC \#\#\#}}
\index{PET 2001}
\item CBM 4032, CBM 8032, CBM 8096, CBM 8296 
\item CBM 610, CBM 710
\end{itemize}
Commodore Amiga und PC gehören nicht dazu, allerdings existieren 
für diese Rechner andere Lösungen, wie etwa der HxC Floppy Emulator.

\subsection{Über \glqq XD-2031\grqq}
Der Name steht in der Tradition der Namensgebung einer ganzen Familie
von Adapterkabeln, genannt die X-Kabel. Diese konnten zum 
Datenaustausch zwischen dem Diskettenlaufwerk \mbox{VC 1541} 
und PC verwendet werden, als diese noch über eine
Druckerschnittstelle nach Centronics-Standard verfügten.

Das X symbolisiert den Datenaustausch, D steht für \glqq drive\grqq\ 
und 2031 erinnert an das Diskettenlaufwerk CBM 2031.

Auch wenn die Namensgebung die Emulation eines bestimmten Laufwerktyps
vermuten lassen könnte, ist dies nicht der Fall. Es wird auch keine
in den Original-Laufwerken enthaltene 6502 CPU emuliert.
Tatsächlich werden auf DOS-Ebene Dateianfragen beantwortet. Die
Verwendung von Disk-Images (D64, D80 ...) 
ist allerdings möglich.

\subsection{Unterstützte Adapter}
\index{Adapter}
\index{Ger\"at}
Derzeit läuft die Firmware auf folgenden Geräten:

\begin{tabular}[c]{l l}
\toprule 
Hardware & Homepage \\
\midrule
XS-1541	& \url{http://xs1541.t-winkler.net} \\
Raspberry Pi mit 3.1541 & -- \\
petSD	& \url{http://petsd.net} \\
\bottomrule
\end{tabular}
\index{XS-1541!Homepage}
\index{3.1541!Homepage}
\index{petSD!Homepage}

Für nähere Beschreibungen dieser Geräte sei auf deren Internetseiten
verwiesen.

\subsection*{Leistungsmerkmale nach Adapter}
\index{Leistungsmerkmale der Adapter}
\index{Features der Adapter}

\begin{tabular}[c]{l l l l}
\toprule
			& XS-1541	& 3.1541       & petSD		\\
\midrule
IEC-Bus			& $\checkmark$	& $\checkmark$ &  --		\\
IEEE-488-Bus		& $\checkmark$	& --	       & $\checkmark$	\\
Zugriff auf XD-2031 Serverdienste
			& $\checkmark$	& $\checkmark$ & $\checkmark$	\\
SD-Karte (FAT)		& --		& --	       & $\checkmark$	\\
\bottomrule
\end{tabular}

\subsection{Unterstützte Betriebssysteme}
\index{Betriebssysteme}
\begin{tabular}[c]{l c c c}
\toprule
& Linux & OS X & Windows \\
fs-Provider (Dateisystem) & $\checkmark$ & $\checkmark$ & $\checkmark$ \\
di-Provider (Disk Images) & $\checkmark$ & $\checkmark$ & $\checkmark$ \\
ftp-Provider (nur lesend) & $\checkmark$ & $\checkmark$ & $\checkmark$ \\
http-Provider (nur lesend) & $\checkmark$ & $\checkmark$ & $\checkmark$ \\
tcp-Provider (Telnet) & $\checkmark$ & $\checkmark$ & -- \\
Automatischer Serverstart & $\checkmark$ & -- & -- \\
\bottomrule
\end{tabular}

\section{Installation}
\subsection{USB-Treiber}
\index{USB-Treiber}
\index{Treiber}
\index{FT232}
Unter Linux ist der Treiber schon seit einiger Zeit direkt im Kernel
enthalten, so dass keine separate Installation erfolgen muss.

Für OS X und Windows laden Sie sich bitte den Treiber von der
Herstellerseite \url{http://www.ftdichip.com/Drivers/VCP.htm} 
herunter und installieren ihn wie auf Ihrem System üblich.

\subsection{Raspberry Pi unter Raspbian mit 3.1541}
\index{Raspberry Pi!Server einrichten}
\index{3.1541!Server einrichten}
\index{Server!Installation!Raspberry Pi mit 3.1541}
Das 3.1541 wird direkt an den UART des Raspberry Pi angeschlossen.
Die gängigen Linux-Systeme verwenden diesen UART allerdings von
Hause aus als Konsole um unter anderem Kernel-Meldungen auszugeben. 
Da diese Meldungen die Datenübertragung zwischen Raspberry Pi
und dem Coprozessor stören würden, muss der UART erst frei
gegeben werden. Damit der R-Pi den Coprozessor zurücksetzen kann,
wird die Wiring Pi Bibliothek benötigt.

\subsubsection*{Kernelmeldungen nicht mehr auf UART ausgeben}
Die Datei \texttt{/boot/cmdline.txt} enthält die Kernelparameter, die ihm
beim Start übergeben werden.  
Erstellen Sie eine Sicherungskopie, nur für den Fall der Fälle:

\begin{verbatim}
sudo cp /boot/cmdline.txt /boot/cmdline.txt~
\end{verbatim}

Bearbeiten Sie nun mit Root-Rechten die Datei \texttt{/boot/cmdline.txt}:

\begin{verbatim}
sudo vi /boot/cmdline.txt
\end{verbatim}

Entfernen Sie dort alle Parameter, die den seriellen Port \texttt{ttyAMA0}
ansprechen, was am Beispiel hier die folgenden waren:

\begin{verbatim}
console=ttyAMA0,115200 kgdboc=ttyAMA0,115200
\end{verbatim}

Die verbleibende Zeile lautete hier:
\begin{verbatim}
dwc_otg.lpm_enable=0 console=tty1 root=/dev/mmcblk0p2 
rootfstype=ext4 elevator=deadline rootwait
\end{verbatim}

Speichern Sie Ihre Änderungen.

\subsubsection*{Kein Login mehr über UART}
Bearbeiten Sie nun wieder unter
Root-Rechten die Datei \texttt{/etc/inittab}:

\begin{verbatim}
sudo vi /etc/inittab
\end{verbatim}

Durchsuchen Sie die Datei nach \glqq{}ttyAMA0 \grqq{}, bis Sie 
etwa folgendes haben:

\begin{verbatim}
T0:23:respawn:/sbin/getty -L ttyAMA0 115200 vt100
\end{verbatim}

Dieser Eintrag erlaubt es, sich über die serielle Schnittstelle
\texttt{/dev/ttyAMA0} mittels eines Terminals anzumelden.
Damit der Datenverkehr vom 3.1541 nicht als Benutzername und Passwort
missverstanden wird, wird dieser Eintrag durch das Voranstellen einer Raute
auskommentiert:

\begin{verbatim}
#T0:23:respawn:/sbin/getty -L ttyAMA0 115200 vt100
\end{verbatim}

Speichern Sie Ihre Änderungen und starten Sie den Raspberry Pi neu.

\begin{verbatim}
sudo shutdown -r now
\end{verbatim}

\subsubsection*{Installation von Wiring Pi}
\index{Wiring Pi}
\index{Raspberry Pi!Wiring Pi}
Falls bei Ihrer Linux-Distribution \textit{Wiring Pi}\footnote{\url{http://wiringpi.com/}} 
noch nicht enthalten ist, muss es noch installiert werden:

\begin{verbatim}
sudo apt-get update
sudo apt-get upgrade
sudo apt-get install git-core
git clone git://git.drogon.net/wiringPi
cd wiringPi
./build
\end{verbatim}

Die weitere Installation erfolgt wie bei jedem anderen Debian/Ubuntu-artigem Linux.

\label{serverinst}
\subsection{Installation der Server-Software (Linux)}
\index{Server!Installation!Linux}
Für Linux stehen leider noch keine Binärpakete zur Verfügung,
so dass Sie hier aus dem Quelltext selbst compilieren müssen.
Das ist allerdings leichter, als es klingt.

\subsubsection{Voraussetzungen}
Bevor der Server aufgespielt werden kann, müssen erst alle
benötigten Programmpakete installiert werden. Dieser Schritt ist nur
einmalig erforderlich und entfällt bei späteren Aktualisierungen
der Server-Software.

Alle für die weitere Installation benötigten Pakete sollten
sich durch einen Klick auf diesen 
\href{apt://build-essential,gcc-avr,avr-libc,avrdude,avrdude-doc,libcurl4-gnutls-dev,git}{apturl-Link} installieren lassen.

Ansonsten kopieren Sie folgende Anweisung in die Kommandozeile, 
achten Sie dabei darauf, dass alles in \textit{einer} Zeile steht:

\begin{verbatim}
sudo apt-get install build-essential gcc-avr avr-libc 
avrdude avrdude-doc libcurl4-gnutls-dev git
\end{verbatim}

Sie haben nun die Wahl, ob Sie den Entwicklerstand oder eine
bestimme Release-Version verwenden möchten. Ein Entwicklerstand
enthält die aktuelle Fassung und enthält die neuesten Features,
Fehlerbehebungen und möglicherweise auch neue Fehler.

\subsubsection{Installation aus dem Entwicklerstand}
\index{Quelltext}
\index{Entwicklerstand}
\index{github.com}

\begin{verbatim}
git clone https://github.com/fachat/XD2031
cd XD2031/pcserver
make
sudo make install
\end{verbatim}

Fertig. Wenn Sie später eine neuere Fassung verwenden möchten, müssen Sie
nicht wieder alles herunter laden, sondern können folgende Befehle
anwenden:

\begin{verbatim}
cd XD2031/pcserver
git pull
make
sudo make install
\end{verbatim}

\subsubsection{Installation einer Release-Version}
Laden Sie von \srcdownloads{}  den Quelltext her"-unter und
\index{Quelltext}
entpacken Sie ihn. Wechseln Sie in das Quelltextverzeichnis
und führen die folgenden Befehle aus:

\begin{lstlisting}
cd XD2031/pcserver
make
sudo make install
\end{lstlisting}

\subsubsection{Automatischer Start des Servers}
\index{Automatischer Serverstart!Linux}
Es bietet sich an, den Server beim Hochfahren automatisch
starten zu lassen. Ihre zugewiesenen Inhalte stehen dann sofort
bereit, wenn Sie sie brauchen. Da der Server nur aktiv wird,
wenn Daten ausgetauscht werden, kostet Sie das bei Nichtbenutzung praktisch
keine Rechenleistung.

Bei Geräten, die mittels USB angeschlossen werden,
können Sie diese auch ab- und anstecken, wobei der Server direkt bereit sein
wird. Bei manuellem Start ist dies nicht möglich, ein Abstecken wird den
Server beenden.

Der verwendete Autostart-Mechanismus basiert auf den klassischen SysV-Style,
sollte aber auch von neueren Systemen unterstützt werden, die systemd oder
Upstart verwenden.

Wenn Sie den Server beim Hochfahren des Rechners automatisch mit starten
soll, führen Sie nach \texttt{sudo make install} noch den folgenden Befehl aus:

\begin{verbatim}
sudo make daemon
\end{verbatim}

Bei der Installation wird eine Beispiel-Konfiguration abgelegt, die Sie
\index{Server!Konfiguration}
mit Root-Rechten nach Ihrem Gusto bearbeiten können, wozu Sie bitte
Einzelheiten zu den Parametern des Server-Starts ab Seite \pageref{serverconf}
dem Kapitel \glqq\ref{serverconf} \nameref{serverconf}\grqq\
entnehmen. Nach jeder Änderung muss der Server neu gestartet werden,
damit er der Änderungen gewahr wird.


\begin{verbatim}
sudo vi /etc/default/fsser
service fsser restart
\end{verbatim}

Meldungen des Servers werden in \texttt{/var/log/fsser}
abgelegt. Sie können wie folgt live mitgelesen werden:
\index{Server!Meldungen}
\index{/var/log/fsser}

\begin{verbatim}
tail -f /var/log/fsser
\end{verbatim}

Den Rest dieses Kapitels werden Sie wahrscheinlich überspringen können
und nur bei Bedarf zu Rate ziehen.

\clearpage
\subsubsection*{Init-Skript, Daemon und Server}
\index{Init-Skript}
\index{/etc/default/fsser}
\index{/usr/local/bin/fsser}
\index{Daemon}
\index{service (Befehl)}
Der selbst-startende Server besteht aus insgesamt fünf Elementen:
\begin{enumerate}
\item Konfigurations-Datei \texttt{/etc/default/fsser}
\item Init-Skript \texttt{/etc/init.d/fsser}
\item Daemon \texttt{/usr/local/bin/fsserd}
\item Server \texttt{/usr/local/bin/fsser}
\item Log-Datei \texttt{/var/log/fsser}
\end{enumerate}

Das Init-Skript wird beim Systemstart automatisch mit Root-Rechten
aufgerufen. Zunächst liest die Konfigurations-Datei, die später
auch vom Daemon verwendet wird. Es stellt sicher, dass
die Log-Datei \texttt{/var/log/fsser} verfügbar ist und durch
den zugewiesenen Benutzer beschreibbar ist. Zuletzt wird der
Daemon unter den Rechten des zugewiesenen Benutzers gestartet.

Der Daemon bleibt aktiv und überwacht regelmäßig die Verfügbarkeit der zugewiesenen
Geräteschnittstelle. Sobald diese verfügbar wird, startet der Daemon einen
Server, der die eigentliche Arbeit leistet und die Dateien für
den CBM bereit stellt.
Weiterhin überwacht der Daemon, ob der Server noch läuft.
Falls dies nicht mehr der Fall sein sollte, prüft der Daemon nach gewissen Regeln,
ob und wann der Server neu gestartet werden soll.

\subsubsection*{Konfiguration des Init-Skripts}
Über die Datei \texttt{/etc/default/fsser} wird das Init-Skript
konfiguriert. Zum Bearbeiten dieser Datei sind Root-Rechte erforderlich.
Nach jeder Bearbeitung muss der Server neu gestartet werden, damit die
aktualisierte Datei eingelesen wird:

\begin{verbatim}
sudo vi /etc/default/fsser
service fsser restart
\end{verbatim}

\texttt{DAEMON\_ARGS} \index{DAEMON\_ARGS} sind beinahe beliebige Parameter,
die dem Server beim Start übergeben werden.
Tragen Sie hier Ihre ASSIGN-Befehle ein um festzulegen, welche Inhalte der
Server bereit stellen soll.
Übergeben Sie hier keine -d Parameter zur Angabe der Gerätedatei,
sondern verwenden Sie statt dessen \texttt{DEVICE}.

\texttt{DEVICE} \index{DEVICE} definiert die zu überwachende Gerätedatei, die für den
Server verwendet werden soll. Sie können Sie in den meisten Fällen undefiniert
lassen, was die Verwendung des Vorgabewerts \texttt{auto} bewirken wird.
Für die meisten Fälle wird dann \texttt{/dev/ttyUSB0} verwendet werden.
Auf dem Raspberry Pi wird diese Schnittstelle ebenfalls verwendet, wenn das
Gerät beim Systemstart bereits verfügbar ist, andernfalls wird die
fest eingebaute serielle Schnittstelle \texttt{/dev/ttyAMA0} verwendet, die
beispielsweise vom 3.1541 verwendet wird.

\texttt{RESPAWN}\ \index{RESPAWN} legt fest, ob im Fehlerfall ein Neustart des Servers
gewünscht ist (\texttt{=1}) oder nicht (alle anderen Werte,
z.B. \texttt{RESPAWN=0}). Dies ist keine starre Regel, sondern eher ein
Anhaltspunkt. Wenn beispielsweise das Gerät abgesteckt wird, beendet das
zwangsläufig den Server. Dies wird nicht als echter Fehler behandelt,
sondern der Server wird immer neu gestartet werden, wenn das Gerät wieder
verfügbar ist. Bei schweren, nicht behebbaren Fehlern, wie beispielsweise
fehlerhafte Konfigurationsparametern wird niemals neu gestartet.
Diese Option behandelt eher den unvorhergesehen Fall und lässt Ihnen die
Wahl, ob Ihre Priorität auf der Verfügbarkeit liegt oder auf dem Finden
von Umständen, die den Server außerplanmäßig beenden.

\texttt{RUN\_AS\_USER} \index{RUN\_AS\_USER} legt den Benutzer fest,
unter dem der Server läuft.
Das Init-Skript wird beim Systemstart als root aufgerufen, und es ist
eine sehr gute Idee, den Server nicht ständig mit root-Rechten laufen zu lassen.
Bei der Installation wird Ihr normaler Benutzername (aus der
Umgebungsvariable \texttt{SUDO\_USER} eingetragen. Sie können dem Server
auch einen eigenen Benutzer zuweisen, werden dann aber alle Aufrufe des
Init-Skripts (z.B. zum Neustart des Servers) mittels \texttt{sudo}
durchführen müssen.

Wenn Sie die Parameter in \texttt{/etc/default/fsser} bearbeitet haben,
müssen Sie den Server neu starten, damit er der Änderungen gewahr wird. Dies
wird im folgenden Kapitel beschrieben.

\subsubsection*{Manuelle Aufrufe des Init-Skripts}
\index{Init-Skript}
\index{Server!Neustart}
Neustart des Servers:

\begin{verbatim}
service fsser restart
\end{verbatim}

Abfrage des Status von Server und Daemon:
\index{Server!Status}

\begin{verbatim}
service fsser status
\end{verbatim}

Manueller Start/Stopp von Server und Daemon:

\begin{verbatim}
service fsser start
service fsser stop
\end{verbatim}

\subsubsection*{Automatischen Start deaktivieren}
Wenn Sie den Mechanismus des automatischen Starts entfernen möchten,
und den Quelltext noch zur Hand haben, wenden Sie folgende Befehlszeilen an:

\begin{verbatim}
cd XD2031/pcserver
sudo make exorcism
\end{verbatim}

Manuell kann der Autostart wie folgt entfernt werden:

\begin{verbatim}
service fsser stop
sudo rm /etc/init.d/fsser
sudo rm /usr/local/install/fsserd
sudo rm /etc/default/fsser
sudo update-rc.d fsser remove
\end{verbatim}

Dabei werden zunächst eventuell laufende Server und überwachende Init-Skripte
beendet. Anschließend wird das Init-Skript, der Daemon, die persönliche Konfiguration
und schließlich die Links der verschiedenen Run-Level
auf das Init-Skript entfernt.

\subsection{Installation der Server-Software (OS X)}
\index{Server!Installation!OS X}
Laden Sie von \osxbinaries{} den Server her"-unter und installieren ihn 
wie üblich.

\subsection{Installation der Server-Software (Windows)}
\index{Server!Installation!Windows}
Laden Sie von \winbinaries{} den Server her"-unter. Entpacken Sie die
zip-Datei in einem Verzeichnis Ihrer Wahl. Der Server ist von dort aus
direkt lauffähig, eine gesonderte Installation ist nicht erforderlich.

\subsection{Aktualisierung der Firmware}
\index{Firmware}
Die XD-2031-Software, die auf einem Mikrocontroller-basiertem Adapter
läuft, wird auch als \textit{Firmware} bezeichnet.

Sicherheitshalber sollten Sie nach dem Installieren einer neuen
Server-Version auch die Firmware aktualisieren, damit beide richtig
zusammen arbeiten können.

Anweisungen für Raspberry Pi / 3.1541 finden Sie unter
bei \glqq\ref{fwinst31541} \nameref{fwinst31541}\grqq\ auf Seite
\pageref{fwinst31541}.

Wenn Sie ein petSD besitzen, lesen Sie bitte jetzt bei
\glqq\ref{petsdfwinst} \nameref{petsdfwinst}\grqq\ auf Seite
\pageref{petsdfwinst} weiter.

\subsubsection{Aktualisierung der Firmware für XS-1541}
\index{XS-1541!Firmware aktualisieren}
Laden Sie zunächst von \fwbinaries{} die Firm"-ware für das XS-1541
herunter und entpacken Sie wenn nötig Archivdateien wie 
etwa \glqq .zip\grqq . 

\clearpage
\subsubsection*{Linux}
\index{XS-1541!Firmware aktualisieren!Linux}
\begin{enumerate}
\item Schließen Sie das XS-1541 mittels USB-Kabel an Ihren Computer an.
\item Öffnen Sie ein Terminal.
\item Wechseln Sie in das Verzeichnis, das die entpackte Firmware enthält.
\item Machen Sie die Skript-Datei xs1541up.sh wie folgt ausführbar:

\texttt{chmod +x xs1541up.sh}
\item Für die aktuelle Firmware \texttt{XD2031-xs1541.hex}, tippen Sie
	den folgenden Befehl ein, aber noch \textbf{nicht} die
	Eingabetaste:

	\texttt{./xs1541up.sh}

\item Drücken Sie am XS-1541 die Reset-Taste und lassen Sie sie wieder los.
\item Drücken Sie am PC binnen drei Sekunden(!) die Eingabetaste
\end{enumerate}

\textbf{Tipp:} Wenn Sie statt der vorgegebenen Datei 
\texttt{XD2031-xs1541.hex}
eine andere Firmware aufspielen möchten, können Sie
dem Skript auch einen Dateinamen übergeben:

\begin{verbatim}
./xs1541up.sh Beliebige-xs1541-Firmware.hex
\end{verbatim}

\subsubsection*{OS X}
\index{XS-1541!Firmware aktualisieren!OS X}
\begin{enumerate}
\item Installieren Sie das \textit{CrossPack for 
AVR\textsuperscript{\textregistered} Development} von 
\\\url{http://www.obdev.at/products/crosspack}
\item Schließen Sie das XS-1541 mittels USB-Kabel an Ihren Computer an.
\item Klicken Sie auf \glqq Programme\grqq\ $\to$ \glqq Dienstprogramme\grqq
$\to$ \glqq Terminal\grqq\ um ein Terminal zu öffnen.
\item Wechseln Sie in das Verzeichnis, das die entpackte Firmware enthält.
\item Machen Sie die Skript-Datei xs1541up.sh wie folgt ausführbar:

\texttt{chmod +x xs1541up.sh}
\item Tippen Sie den folgenden Befehl ein, aber noch \textbf{nicht} die
Eingabetaste:

\texttt{sudo mv /dev/cu.usbserial}

\item Tippen Sie jetzt die Tab-Taste auf Ihrer Tastatur. Das Wort
\glqq usbserial\grqq\ sollte jetzt mit der Seriennummer Ihres Adapters
vervollständigt werden, wie z.B. usbserial-A10044Co.

Drücken Sie noch immer \textbf{nicht} die Eingabetaste, sondern fügen
Sie die folgende Zeile an den Befehl an:

\textit{<Leertaste>} \texttt{/dev/ttyUSB0}

Wenn der Befehl jetzt wie folgt aussieht (ausschließlich anderer
Seriennummer), drücken Sie jetzt die Eingabetaste:

\texttt{sudo mv /dev/cu.usbserial-A10044Co /dev/ttyUSB0}

Geben Sie dann Ihr Administrator-Passwort ein, um diesen Vorgang
zu ermöglichen.
\item Für die aktuelle Firmware 
   \texttt{XD2031-xs1541.hex} tippen Sie
	folgenden Befehl ein, aber noch \textbf{nicht} die
	Eingabetaste:

	\texttt{./xs1541up.sh}

\item Drücken Sie am XS-1541 die Reset-Taste und lassen Sie sie wieder los.
\item Drücken Sie am PC binnen drei Sekunden(!) die Eingabetaste
\end{enumerate}

\textbf{Tipp:} Wenn Sie statt der vorgegebenen Datei 
\texttt{XD2031-xs1541.hex}
eine andere Firmware aufspielen möchten, können Sie
dem Skript auch einen Dateinamen übergeben:

\begin{verbatim}
./xs1541up.sh Beliebige-xs1541-Firmware.hex
\end{verbatim}

\subsubsection*{Wenn das Aktualisieren scheitert...}
Wenn es nicht geklappt hat, weil drei Sekunden doch recht kurz sind,
müssen Sie nicht alles noch einmal eintippen. Drücken Sie einfach
die $\uparrow$-Taste und der letzte Befehl sollte
wieder erscheinen. Drücken Sie dann Reset und jetzt schnell die Eingabetaste!

Sollte trotz allem die Fehlermeldung \glqq programmer is not responding\grqq\ 
erscheinen, wurde Ihr XS-1541 vermutlich \textit{nicht} mit einem Bootloader
versehen. In diesem Fall kann die Firmware nur unter Verwendung eines
AVR-ISP-Programmierers aufgespielt werden.



\subsubsection{Aktualisierung der Firmware für Raspberry Pi / 3.1541}
\index{3.1541!Firmware aktualisieren}
\label{fwinst31541}
Laden Sie zunächst von \fwbinaries{} die Firm\-ware für das 3.1541
herunter und entpacken Sie wenn nötig Archivdateien wie etwa \glqq .zip\grqq . 
Das 3.1541 verwendet die selbe Firmware die das XS-1541, sie werden daher
keine separaten Dateien finden. Allerdings ist der Update-Vorgang etwas
leichter:

\begin{enumerate}
\item Wechseln Sie in das Verzeichnis, das die entpackte Firmware enthält.
\item Machen Sie die Skript-Datei 3.1541.sh wie folgt ausführbar:

\texttt{chmod +x 3.1541up.sh}
\item Um die aktuelle Firmware \texttt{XD2031-xs1541.hex} aufzuspielen:

	\texttt{./3.1541up.sh}
\end{enumerate}

\textbf{Tipp:} Wenn Sie statt der vorgegebenen Datei 
\texttt{XD2031-xs1541.hex}
eine andere Firmware aufspielen möchten, können Sie
dem Skript auch einen Dateinamen übergeben:

\begin{verbatim}
./3.1541up.sh Beliebige-xs1541-Firmware.hex
\end{verbatim}

\subsubsection{Aktualisierung der Firmware für petSD}
\index{petSD!Firmware aktualisieren}
\label{petsdfwinst}
Zwar wurden die 2011er petSD mit einem Bootloader zum einfachen
Aktualisieren der Firmware über die SD-Karte versehen, es hat
sich jedoch gezeigt, dass die verwendeten Fuses-Einstellungen
des AVR Mikrocontroller nicht mit der
\mbox{XD-2031} Firmware kompatibel sind.

Leider können Bootloader nur die Firmware aktualisieren, 
aber keine Änderungen an den Fuses vornehmen. 
Die einzige Möglichkeit, die XD-2031 Firmware
auf ein petSD aufzuspielen, besteht daher in der Verwendung eines
AVR-ISP-Programmiergerätes. 

Ein solches Programmiergerät kann je nach Hersteller 
für einen Betrag von ca. 20,-- bis 30,-- \euro{} erworben werden. 
Das Atmel AVR ISP mkII ist eine Empfehlung wert und wird in allen
Skripten als Standard verwendet.

\subsubsection*{Download der Firmware und Einstellen des
Programmier-Gerätes}
\begin{enumerate}
\item Unter OS X installieren Sie bitte erst das 
\textit{CrossPack for AVR\textsuperscript{\textregistered} Development} von 
\url{http://www.obdev.at/products/crosspack}

\item Laden Sie von \fwbinaries{} die Firmware herunter und entpacken Sie 
die Datei

\item Öffnen Sie ein Terminal. Unter OS X klicken Sie dazu auf 
\glqq Programme\grqq\ $\to$ \glqq Dienstprogramme\grqq
$\to$ \glqq Terminal\grqq{}

\item Wechseln Sie in das Verzeichnis, das die entpackte Firmware enthält
\end{enumerate}

Wenn Sie ein anderes Programmiergerät als das Atmel AVR ISP mkII
verwenden, müssen die zunächst die für Ihr Programmiergerät passenden
Werte eintragen. Lesen Sie dazu die Anleitung Ihres Programmiergerätes oder 
die Online-Hilfe mit \texttt{man avrdude}

\begin{enumerate}
\setcounter{enumi}{4}
\item Bearbeiten Sie wenn nötig die Datei \texttt{programmer.cfg} 
und ändern Sie die beiden Parameter für Typ und Anschluss des 
Programmiergerätes in pas\-sende Werte:
\begin{verbatim}
PROGRAMMER=avrispmkii
CONNECTED_TO=usb
\end{verbatim}
\end{enumerate}

\subsubsection*{Anpassen der Fuses}
Dieser Schritt ist nur einmalig notwendig und sollte bei
weiteren Firmware-Updates ausgelassen werden.
\begin{enumerate}
\item Machen Sie das Skript \texttt{petSD\_fix\_fuses.sh} ausführbar:
\begin{verbatim}
chmod +x petSD_fix_fuses.sh
\end{verbatim}

\item Starten Sie das Skript, um die Fuses für XD-2031 zu ändern:
\begin{verbatim}
./petSD_fix_fuses.sh
\end{verbatim}

\end{enumerate}

\subsubsection*{Aktualisieren der Firmware mit AVR ISP Programmierer}

\begin{enumerate}
\item Machen Sie das Skript \texttt{petSDisp.sh} ausführbar:
\begin{verbatim}
chmod +x petSDisp.sh
\end{verbatim}

\item Um die aktuelle Firmware \texttt{XD2031-petSD.hex} aufzuspielen:

\texttt{./petSDisp.sh}
\end{enumerate}

\textbf{Tipp:} Wenn Sie statt der vorgegebenen Datei 
\texttt{XD2031-petSD.hex}
eine andere Firmware aufspielen möchten, können Sie
dem Skript auch einen Dateinamen übergeben:

\begin{verbatim}
./petSDisp.sh Beliebige-petSD-Firmware.hex
\end{verbatim}

\begin{comment}
Haben Sie Ihr petSD nach Juli 2012 erworben, können Sie die Firmware
sehr einfach aktualisieren:
\begin{enumerate}
\item Laden Sie von \fwbinaries{} die Firmware für das \mbox{petSD}
herunter und entpacken Sie wenn nötig Archivdateien wie etwa \glqq .zip\grqq . 
\item Kopieren Sie die auf \glqq .bin\grqq\ endende Firmware-Datei
in das Hauptverzeichnis der am petSD verwendeten SD-Karte
\item Legen Sie die SD-Karte in das ausgeschaltete petSD ein
\item Schalten Sie das petSD ein
\end{enumerate}
Der Bootloader wird dann die neue Firmware finden und den Mikrocontroller
neu programmieren. Die Datei darf auch nach dem Programmiervorgang im 
Hauptverzeichnis liegen bleiben, da der Bootloader gleiche Firmwares
nur einmal programmiert.

Haben Sie dagegen Ihr petSD vor Juli 2012 erworben, kann es zu
Problemen kommen. Schreiben Sie in diesem Fall an nils.eilers@gmx.de, 
um weitere Informationen zu erhalten.
\end{comment}

\clearpage
\section{Erste Schritte}
Die meiste Arbeit ist nun geschafft. Bevor wir nun zum vergnüglichen
Teil kommen, sollen kurz drei Begriffe geklärt werden, die von nun
an in dieser Anleitung verwendet werden:
\begin{itemize}
\item Ihr moderner Computer wird fortan der Einfachheit halber
	\glqq{}PC\grqq{} genannt. \index{PC}
\item Ihr alter Commodore wird als 
\glqq{}CBM\grqq{} bezeichnet, ganz gleich, ob es sich um einen CBM / PET / C64 
oder anderen 8-Bitter handelt.
\index{CBM} \index{C64} \index{PET 2001} \index{C128}
\item Ihr XS-1541, 3.1541 oder petSD wird 
als \glqq{}Adapter\grqq{} bezeichnet.
\end{itemize}

\subsection{Adapter anschließen}
\index{Anschließen des Adapters}
Auch wenn es hier erwähnt wird, sollten Sie vorerst auf das Anschließen
weiterer Diskettenlaufwerke verzichten, um Konflikte um Geräteadresse 8
auszuschließen.
\begin{enumerate}
\item Stellen Sie sicher, dass am Adapter \textit{keine(!)} Kabel angeschlossen sind.
\item Stellen Sie sicher, dass der CBM ausgeschaltet ist.
\item Wenn neben dem CBM noch weitere Hardware am IEEE-488-Bus oder IEC-Bus
angeschlossen ist, wie etwa weitere Diskettenlaufwerke oder Drucker,
müssen diese ebenfalls alle ausgeschaltet sein.
\item Der PC darf zu diesem Zeitpunkt bereits eingeschaltet sein. Falls er das noch nicht ist, schalten Sie ihn jetzt ein.
\item Verbinden Sie den Adapter mit dem CBM über ein IEEE-488- oder IEC-Bus-Kabel.
\item Verbinden Sie den Adapter mit dem PC mittels USB-Kabel.
\item Schalten Sie Peripherie wie Diskettenlaufwerke oder Drucker ein.
\item Wenn der Adapter über eine separate Stromzufuhr verfügt, schalten Sie diese jetzt ein.
\item Schalten Sie den CBM ein.
\end{enumerate}

Vielleicht staunen Sie über dieses Prozedere. Tatsache ist, das weder
IEEE-488-Bus noch IEC-Bus \textit{hot-pluggable} sind; das bedeutet:
Ein- oder Abstecken nur im ausgeschalteten Zustand -- auch wenn man es in
der Praxis immer wieder anders sieht. Nichtbeachtung kann erforderlich
machen, dass Sie Ersatzteile für Ihren CBM besorgen müssen und Chips
aus- und einlöten müssen. Nebenbei bemerkt: der Bustreiber MC3446 ist
teuer geworden.

Falls Sie ein XS-1541 verwenden, und doch erst den CBM einschalten, bevor
das XS-1541 eingeschaltet ist, werden Sie ein Glimmen der Power-LED
beobachten können. Das rührt daher, dass die Spannung, die am Bus auch im
Ruhezustand anliegt, über die internen Sicherungsdioden des Mikrocontrollers
in das Stromnetz des Adapters fließt. Das ist ein Zustand, der
eindeutig außerhalb der \glqq absolute maximum ratings\grqq\ liegt,
die der Hersteller Atmel in seinem Datenblatt nennt. Atmel warnt,
dass dadurch die Zuverlässigkeit beeinträchtigt werden kann und bleibende
Schäden entstehen können. 

\subsection{XS-1541 ein- und ausschalten}
\index{Einschaltreihenfolge}
Natürlich wollen wir den Adapter jetzt eingeschaltet lassen, aber dies
ist die richtige Stelle, die Ein- und Ausschaltreihenfolge für das
XS-1541 zu besprechen. Beim petSD ist die Reihenfolge unerheblich.

\subsection*{XS-1541 einschalten}
\begin{enumerate}
\item Schalten Sie immer erst den PC ein. 
\item Stellen Sie sicher, dass das USB-Kabel am XS-1541 eingesteckt ist und 
die Power-LED hell leuchtet.
\item Wenn vorhanden, schalten Sie jetzt weitere Peripherie wie 
Diskettenlaufwerke oder Drucker ein.
\item Zuletzt schalten Sie den CBM ein, damit alle angeschlossenen
Geräte in den Grundzustand versetzt werden.
\end{enumerate}

\subsection*{XS-1541 ausschalten}
\index{Ausschalten, Reihenfolge}
\begin{enumerate}
\item Schalten Sie den CBM und alle am Bus hängenden Peripheriegeräte
wie Diskettenlaufwerke und Drucker aus.
\item Ziehen Sie erst jetzt den USB-Stecker vom XS-1541 ab bzw.
schalten Sie erst jetzt den PC aus.
\end{enumerate}

\subsection{Server starten}
\index{Server!Starten}
Die Konfiguration des Servers ist für diese Beispiele nicht erforderlich.
Sie wird später beschrieben.
Öffnen Sie ein Terminal. Unter OS X klicken Sie dazu auf \glqq Programme\grqq\  
$\to$ \glqq Dienstprogramme\grqq $\to$ \glqq Terminal\grqq\ .

Starten Sie dann den Server mit seiner Standard-Konfiguration, indem Sie
den folgenden Befehl eingeben und mit der Eingabetaste bestätigen:

\begin{verbatim}
fsser
\end{verbatim}


\subsection{Inhaltsverzeichnis laden}
\index{Inhaltsverzeichnis}
Gegenüber dem CBM erscheint der Adapter wie ein klassisches Diskettenlaufwerk
unter einer einzigen Geräteadresse. Voreingestellt ist Geräteadresse 8.
So können Sie sich direkt wie früher das Inhaltsverzeichnis anzeigen
lassen:

\begin{verbatim}
LOAD"$",8
LIST
\end{verbatim}

Probieren Sie es aus. Sie sollten das Heimatverzeichnis Ihres PCs
angezeigt bekommen.

\subsection{LOAD und SAVE}
\index{LOAD}
\index{SAVE}
Geben Sie ein kleines Programm ein:

\lstset{language=,} %[Visual]Basic}
\begin{lstlisting}
10 FOR I=1 TO 10
20 PRINT I,I*I
30 NEXT

RUN
\end{lstlisting}

Läuft es? Wenn Sie mit der Ausgabe zufrieden sind, speichern Sie Ihr Werk:

\texttt{SAVE"TEST",8}

Laden Sie dann erneut das Inhaltsverzeichnis:

\begin{verbatim}
LOAD"$",8
LIST
\end{verbatim}

Dort sollte nun neben den bereits bekannten Dateien auch Ihr
kleines Programm erscheinen. Nun hat das Inhaltsverzeichnis 
aber Ihr Programm im Speicher überschrieben, wovon Sie sich
gerne mit \texttt{LIST} oder \texttt{RUN} überzeugen können.
Holen Sie sich Ihr Programm zurück:

\begin{verbatim}
LOAD"TEST",8
RUN
\end{verbatim}


\subsection{Mehrere Laufwerke unter einer Geräteadresse}
\index{Laufwerke}
\index{Ger\"ateadresse}
Unter einer Geräteadresse können mehrere Laufwerke gleichzeitig zur
Verfügung gestellt werden. 

Während das für Benutzer von Doppellaufwerken
wie CBM 8050 selbstverständlich ist, mag dieses Konzept für Benutzer
von Einzellaufwerken wie VC 1541 am C64 ungewohnt sein, weshalb es
hier kurz erklärt werden soll: vor Dateinamen kann die Nummer des gewünschten
Laufwerks durch Doppelpunkt getrennt angegeben werden. Geräte mit
mehreren Laufwerken wissen dann, auf welches Laufwerk der Zugriff
erfolgen soll.

Sehen wir uns doch einfach mal die dämliche Demo auf Laufwerk 1 an:

\begin{verbatim}
LOAD"$1:DUMB*",8
LIST

 1 "DUMB*           "XD2031
 1     "DUMBDEMO64"      PRG
 1     "DUMBDEMOPET"     PRG
 1     "DUMBDEMOVIC20"   PRG
 63999 BLOCKS FREE.
READY.
\end{verbatim}

Nicht gerade viel versprechend, aber sehen Sie selbst... 
Wählen Sie dazu den Dateinamen, der für Ihr System passt.

\begin{verbatim}
LOAD"DUMBDEMOPET",8
RUN
\end{verbatim}

Wieso wurde das Programm geladen, obwohl wir doch \texttt{LOAD"DUMBDEMOPET",8} 
geschrieben haben? Hätten wir nicht \texttt{LOAD"1:DUMBDEMOPET",8} schreiben müssen,
um von Laufwerk 1 zu laden?

\index{Laufwerk!zuletzt verwendetes}
\index{Laufwerk!aktuelles}
Nun, \textit{ohne} Laufwerksnummer wird die \textit{zuletzt verwendete}
Laufwerksnummer verwendet. Und das war die von Laufwerk 1, von dem wir
das Inhaltsverzeichnis geladen haben.

Wenn Sie nun Ihr Testprogramm zurück erhalten wollen, können Sie
es sich mit \texttt{LOAD"0:TEST",8} wieder zurück holen.

XD-2031 erlaubt die Verwendung von bis zu 10 Laufwerken gleichzeitig,
die mit den vorgestellten Laufwerksziffern 0 bis 9 angesprochen
werden.

\subsection{BASIC 4.0 Disketten-Befehle}
\index{BASIC 4.0}
Prinzipiell sind die zusätzlichen Disketten-Befehle des BASIC 4.0
nützlich und werden auch von der Firmware unterstützt. 
Eine gute Gelegenheit, das Hauptverzeichnis Ihres PCs
wieder zu säubern:

\begin{verbatim}
SCRATCH"TEST",D0

ARE YOU SURE ?Y
01,FILES SCRATCHED,01,00
\end{verbatim}

Leider wird die Laufwerksangabe streng geprüft. Bei Angaben größer 1
wird ein \texttt{ILLEGAL QUANTITY ERROR} gemeldet. Dreifachlaufwerke
waren Commodore dann wohl doch nicht geheuer.

Immerhin sind bei der Angabe der Geräteadresse über ONU alle
zulässigen Werte möglich. Mit \texttt{CATALOG ONU 20} würde das
Inhaltsverzeichnis eines Laufwerks an Geräteadresse 20 angezeigt werden.
Werte größer als 31 oder kleiner als drei werden verweigert. 

BASIC prüft hier zwar falsch, da die größtmögliche Geräteadresse nicht 31
sondern nur 30 ist, aber wenigstens stellt das keine Einschränkung dar.

\subsection{Geräteadresse ändern}
\index{Ger\"ateadresse!\"andern}
\index{Adresse|see{Ger\"ateadresse}}
Das Ändern der Geräteadresse Ihres Adapters erfolgt durch Senden
eines Befehls an den Kommandokanal mit Sekundäradresse 15. Eine einfache
Möglichkeit, mit BASIC-Bordmitteln Befehle zu senden, ist das Anhängen
des Befehls an den OPEN-Befehl.

\begin{verbatim}
OPEN 1,8,15,"XU=30" : CLOSE 1
\end{verbatim}

Dies sendet dem Adapter an Geräteadresse 8 über den Kommandokanal 15
den Befehl, fortan die neue Geräteadresse 30 zu verwenden.

Wenn Sie BASIC 4 haben, probieren Sie ein \texttt{CATALOG ONU 30}, ansonsten:
\begin{verbatim}
LOAD"$",30
LIST
\end{verbatim}

Für die Assembler- oder C-Programmierer kann es einfacher sein,
die Laufwerksangabe binär statt ASCII-dezimal zu übertragen. Mit diesem
zugegebenermaßen BASIC-Beispiel hierfür stellen wir die Geräteadresse
wieder von 30 zurück auf 8:

\begin{verbatim}
OPEN 1,30,15
PRINT#1,"XU=";CHR$(8)
CLOSE 1
\end{verbatim}

Nach dem Ausschalten verwendet der Adapter wieder die voreingestellte
Geräteadresse 8. Folgende Software-Versionen werden die Möglichkeit 
bieten, den eingestellten Wert dauerhaft auch über das Ausschalten hinweg
beizubehalten.

\clearpage %MANFF
\subsection{Laufwerken Inhalte zuweisen}
\index{Provider}
Bei klassischen Diskettenlaufwerken würden Sie eine Diskette 
einlegen. Stattdessen muss hier dem Adapter mitgeteilt werden, 
\textit{wo} wo die Daten liegen, die als virtuelle Diskette 
erscheinen sollen.

Weil der Zugriff auf eine Festplatte aber grundsätzlich anders abläuft
als beispielsweise der Zugriff auf einen FTP-Server im Internet, 
muss man dem Adapter auch noch mitteilen \textit{wie} er darauf
zugreifen soll.

Für jeden Speicherort (Dateisystem/FTP/HTTP) gibt es daher einen 
\textit{Provider.} Auf Deutsch bedeutet das etwa \textit{Lieferant}
oder \textit{Leistungserbringer}. 

Hier ein Beispiel mit vier zugewiesenen Laufwerken:

\begin{tabular}[c]{c l l}
\toprule
Laufwerk & Provider & Ort \\
\midrule
0 & Dateisystem & sample/ \\
1 & Dateisystem & tools/ \\
2 & FTP & ftp.zimmers.net \\
3 & HTTP & www.gutenberg.org/ebooks/ \\
\bottomrule
\end{tabular}

Beispiele für jeden Provider finden Sie im Kapitel \ 
\ref{assign} \glqq{}\nameref{assign}\grqq{} ab Seite
\pageref{assign}.

\clearpage
\section{Dateien und Verzeichnisse}
\subsection{Regeln für Datei- und Verzeichnisnamen}
\index{Dateinamen}
\index{Verzeichnisnamen}
Als maximale Länge eines Datei- oder Verzeichnisnamens sollten
16 Zeichen nicht überschritten werden, damit alte Programme nicht
durch unerwartet lange Dateinamen Schwierigkeiten bekommen.

Generell sollten Sie Sonderzeichen in Dateinamen unbedingt 
vermeiden. Absolut verboten sind:
\begin{itemize}
\item Doppelpunkte :
\item Schrägstriche /
\item umgekehrte Schrägstriche \textbackslash
\item der Klammeraffe @
\item das Paragraphenzeichen $\S$
\item der Stern *
\item das Fragezeichen ? 
\item das Dollar-Zeichen \$
\item das Komma ,
\item das Gleichheitszeichen =
\end{itemize}
Leerzeichen und Punkte können dagegen problemlos verwendet werden.
Für ein sichtbares Trennzeichen sollten Sie das Minus-Zeichen '-' gegenüber
dem Unterstrich '\_' bevorzugen, da der Unterstrich auf einem CBM schwierig
einzugeben sein kann.

\subsection{Wildcards -- Dateimasken}
\index{Wildcards}
Wildcards funktionieren im CBM-Umfeld anders, als Sie es von DOS,
Windows oder unixoiden Betriebssystemen gewöhnt sein werden, weshalb 
Sie sich wenigstens den nächsten Abschnitt \glqq \nameref{starwc}\grqq\ 
ansehen sollten.

Aber zunächst sei erst einmal sei erklärt, was Wildcards oder Dateimasken 
überhaupt sind: nämlich
eine Art Joker für Dateinamen, die das Arbeiten mit mehreren Dateien
auf einmal sehr erleichtern.

Das Fragezeichen '?' steht dabei für ein beliebiges vorhandenes einzelnes 
Zeichen, das '*' für mehrere beliebige Zeichen oder nichts.

Angenommen, Sie hätten eine Liste von Dateien:

\begin{verbatim}
BUDGET1977
BUDGET1978
BUDGET1979
BUDGET1979.Q1
BUDGET1979.Q2
BUDGET1979.Q3
\end{verbatim}

\texttt{@\$:BUDGET197?} würde dann BUDGET1977, BUDGET1978 und BUDGET 1979
anzeigen, nicht aber BUDGET1979.Q1, BUDGET1979.Q2 und BUDGET1979.Q3.

\texttt{@\$:BUDGET1979?} würde gar keine Datei anzeigen, da es keine
Datei gibt, deren Name nach \glqq BUDGET1979\grqq\ noch genau ein
weiteres Zeichen enthält.

\texttt{@\$:BUDGET1979*} würde alle vier Dateien BUDGET1979 und 
BUDGET1979.Q1, BUDGET1979.Q2 und BUDGET1979.Q3 anzeigen, da bei
einem Stern nicht notwendigerweise ein Zeichen folgen muss.

\subsubsection*{Arbeitsweise der '*'-Wildcard}
\label{starwc}
Im Unterschied zu anderen Betriebssystem erfolgt nach '*' keine
weitere Auswertung. Wenn Sie also versuchen würden, die Daten
der ersten Quartale durch \texttt{*.Q1} zu kopieren, werden Sie
feststellen, dass statt dessen \textit{alle} Dateien kopiert werden.
Probieren Sie es aus, indem Sie sich mit \texttt{@\$:*.Q1} die
passenden Dateien anzeigen lassen.

\subsection{Verzeichnisse}
\index{Verzeichnisnamen!Trennzeichen}
\index{Trennzeichen}
Die Trennzeichen, die verwendet werden können, um einzelne 
Verzeichnisnamen innerhalb eines Pfades zu trennen, hängen vom
jeweiligen Provider und Betriebs"-system ab. 

Machen Sie es sich einfach, indem Sie ausschließlich
den einfachen Schrägstrich '/' verwenden, der immer und 
überall funktioniert.

Richtig: /Spiele/Strategie\\
Besser nicht: \textbackslash Spiele\textbackslash Strategie

Neben absoluten können auch relative Pfadangaben verwendet werden.
Was bedeutet das?

Nehmen wir an, Sie hätten Ihre Dateien 
in Unterverzeichnissen unter anderem wie folgt strukturiert:

\tikzstyle{every node}=[draw=black,thick,anchor=west]
\tikzstyle{selected}=[draw=red,fill=red!30]
\tikzstyle{optional}=[dashed,fill=gray!50]
\begin{tikzpicture}[%
  grow via three points={one child at (0.5,-0.7) and
    two children at (0.5,-0.7) and (0.5,-1.4)},
      edge from parent path={(\tikzparentnode.south) |- (\tikzchildnode.west)}]
\node {Hauptverzeichnis}
  child { node {Spiele} 
    child { node {Autorennen} }
    child { node {Ballern} }
    child { node [selected] {Strategie} }
  }
%  child [missing] {}
%  child [missing] {}
%  child [missing] {}
%  child { node {Texte} 
%    child { node {Arbeit} }
%    child { node {Briefe} }
%    child { node {Tagebuch} }
%  }
;
\end{tikzpicture}

\index{Pfade}
\index{absolute Pfade|see{Pfade}}
\index{relative Pfade|see{Pfade}}
Nehmen wir weiter an, sie befänden Sich im Unterverzeichnis Strategie
und wollten Sich nun in das Ballerspiel-Verzeichnis wechseln.

Absolute Pfadangaben beinhalten immer den kompletten Pfad vom
Hauptverzeichnis an. Mit \texttt{@CD:/SPIELE/BALLERN} würden Sie nun
das Verzeichnis wechseln.

Bei tiefer verschachtelten Verzeichnissen kann das eine Menge Tipp-Arbeit
bedeuten, weshalb sich relative Angaben anbieten. Zwei Punkte hintereinander
bezeichnen dabei immer das übergeordnete Verzeichnis, ohne dass
dieser Name ausdrücklich bekannt sein müsste. \texttt{@CD:../BALLERN}
bringt Sie dann direkt in das gewünschte Verzeichnis.

Mit \texttt{@CD:..} gelangen Sie immer eine Verzeichnisebene höher.

\texttt{@CD:/} bringt Sie in das Hauptverzeichnis.

\clearpage
\section{DOS-Befehle}
\index{DOS-Befehle}
\index{Syntax}
\subsection{Verwendete Notation}
\index{Notation}
\subsection*{Ersetzung durch tatsächlich passende Werte}
In spitzen Klammern und kursiver Schrift werden Werte angegeben, die 
Sie durch für Sie im Einzelfall tatsächlich
zutreffenden Werte ersetzen müssen. Beispiel:

\texttt{LOAD"}\textit{<Dateiname>}\texttt{",}\textit{<Geräteadresse>}

\texttt{LOAD"TEST",8}

\subsection*{Eine von mehreren Möglichkeiten wählen}
Soll aus mehreren Möglichkeiten genau eine Auswahl getroffen werden,
werden die Wahlmöglichkeiten in Klammern angegeben und durch einen
senkrechten Strich getrennt. Beispiel:

\mbox{\Big( Fisch \Big| \normalsize Fleisch \Big) \normalsize} 

Fleisch

\subsection*{Parameter angeben oder auslassen}
Wenn eine Angabe \textit{optional} ist, können Sie sie entweder
angeben oder weg lassen. Optionale Angaben stehen in eckigen Klammern.

\texttt{LOAD"} \Big[ Laufwerk-Nummer \texttt{:} \Big] \textit{<Dateiname>}
\texttt{",}\textit{<Geräteadresse>}

\texttt{LOAD"TEST",8} lädt die Datei TEST vom zuletzt 
verwendeten Laufwerk

\texttt{LOAD"2:TEST",8} lädt die Datei TEST von Laufwerk 2

\subsection*{Optionale Wiederholung}
Manche Befehle fordern mindestens einen Parameter, können aber beliebig
viele weitere\footnote{
	In der Praxis wird die Anzahl durch andere Umstände begrenzt,
	wie etwa die Größe des Zwischenspeichers, der einen Befehl
	aufnehmen soll
} durch Kommata getrennte Parameter annehmen. Die optionale Wiederholung
wird dann durch ein Komma gefolgt von drei Punkten dargestellt:

\mbox{\textit{<Name>} \Big[ \texttt{,} \ldots \ \Big]}

\begin{verbatim}
PETER
PETER,PAUL,MARY
\end{verbatim}

\subsection{I -- Initialize -- Laufwerk initialisieren}
\index{INIT}

\texttt{I} \Big[ \textit{<Laufwerk-Nummer>} \Big]

Dieser Befehl ist nur aus Kompatibilitätsgründen enthalten.
Seine einzige Wirkung ist das Rücksetzen der Statusmeldung.

Beispiele:
\begin{verbatim}
I
I0
\end{verbatim}

\subsection{R -- Rename -- Datei oder Verzeichnis umbenennen}
\index{RENAME}
\index{Umbenennen}
\index{Datei!umbenennen}
Benennt eine Datei oder ein Verzeichnis um. Beachten Sie die Reihenfolge:
der neue Dateiname wird zuerst genannt.

\texttt{R}\Big[ \textit{<Laufwerk-Nummer>} \Big] \texttt{:}
\textit{<Neuer Name>} \texttt{=} \textit{<Alter Name>}

Beispiele:
\begin{verbatim}
@R:REZEPTE PIZZEN=REZEPT PIZZA
@R0:WUENSCHE2013.TXT=WUENSCHE2012.TXT
\end{verbatim}

\subsection{S -- Scratch -- Datei(en) löschen}
\index{SCRATCH}
\index{Loeschen@L\"oschen!Datei}
\index{Datei!l\"oschen}
\texttt{S} \Big[ \textit{<Laufwerk-Nummer>} \Big] \texttt{:}
\textit{<Dateimaske>} \Big[ \texttt{,} \ldots \  \Big]

Löscht die angegebenen Datei(en). Mehrere Dateinamen bzw. Suchmuster
können durch Kommata getrennt angegeben werden.

Beispiele:

\texttt{@S:MIST} löscht die Datei MIST\\
\texttt{@S:STATISTIK198?.SEQ} löscht die Statistiken der Jahre 1980, 1981 
\ldots \ bis 1989 \\
\texttt{@S:MIST,UNSINN,ENTWURF,ENTWURF2} löscht alle genannten Dateien \\
\texttt{@S:*} löscht alle Dateien(!) des Unterverzeichnisses

\subsection{CD / CHDIR -- Change directory -- Unterverzeichnis wechseln}
\index{Verzeichnis!wechseln}
\index{Ordner|see{Verzeichnis}}
\index{Unterverzeichnis|see{Verzeichnis}}
\mbox{\Big( \texttt{CD} \Big| \texttt{CHDIR} \Big) 
\Big[ \textit{<Laufwerk-Nummer>} \Big]
\texttt{:} \textit{<Verzeichnis>}}

Wechselt in das angegebene Unterverzeichnis.

Beispiele:
\begin{verbatim}
@CD:SPIELE
@CD:..
@CD:TEXTE/HAUSAUFGABEN
\end{verbatim}

\subsection{MD / MKDIR -- Make directory}
\index{MD}
\index{MKDIR}
\index{Verzeichnis!anlegen}

\mbox{\Big( \texttt{MD} \Big| \texttt{MKDIR} \Big) 
\Big[ \textit{<Laufwerk-Nummer>} \Big]
\texttt{:} \textit{<Verzeichnis>}}

Erstellt das angegebene Unterverzeichnis. 

Bei Pfadangaben müssen
alle vorgenannten Verzeichnisse bereits existieren. Bei der
Angabe SOFTWARE/SPIELE/AUTORENNEN müssen also die Verzeichnisse
SOFTWARE/SPIELE bereits existieren, damit das Verzeichnis AUTORENNEN 
erstellt werden kann.

Beispiele:
\begin{verbatim}
@MD:SKIZZEN
@MD:TOPGAMES
\end{verbatim}



\subsection{RD / RMDIR -- Remove directory}
\index{RD}
\index{RMDIR}
\index{Verzeichnis!l\"oschen}
\index{Loeschen@L\"oschen!Verzeichnis}

\mbox{\Big( \texttt{RD} \Big| \texttt{RMDIR} \Big) 
\Big[ \textit{<Laufwerk-Nummer>} \Big]
\texttt{:} \textit{<Verzeichnis>}}

Löscht das angegebene Unterverzeichnis, wenn es weder
Dateien noch Unterverzeichnisse enthält.

Beispiel:
\begin{verbatim}
@RD:EXPERIMENT2
\end{verbatim}


% ----------------------------------------------------------------
\label{xcommands}
% ----------------------------------------------------------------

\subsection{XU -- Change unit -- Geräteadresse ändern}
\index{XU}
\index{Ger\"ateadresse!\"andern}


\texttt{XU=}\Big( \textit{<Geräteadresse in ASCII Ziffern>} \Big|
\textit{<Geräteadresse binär>} \Big)

Soll die Einstellung dauerhaft auch nach dem Ausschalten erhalten 
bleiben, muss die Konfiguration durch den Befehl \texttt{XW} gesichert werden.

Beispiele:

\begin{verbatim}
@XU=9
OPEN 1,8,15,"XU=9" : CLOSE 1
OPEN 1,8,15 : PRINT#1,"XU=";CHR$(9) : CLOSE 1
\end{verbatim}

\subsection{XD - Change last drive - Letztes Laufwerk setzen}
\index{XD}
\index{Laufwerk!zuletzt verwendetes}

\texttt{XD=}\Big( \textit{<Laufwerk-Nummer in ASCII Ziffern>} \Big|
\textit{<Laufwerk-Nummer binär>} \Big)

Das zuletzt verwendete Laufwerk wird als \glqq Standardlaufwerk\grqq\  für
Zugriffe verwendet, bei denen im Dateinamen durch vorangestelltes
\textit{<Laufwerk-Nummer>}: nicht ausdrücklich ein bestimmtes
Laufwerk angegeben wird.  
Zusätzlich kann dieses Laufwerk auch mit dem XD Befehl gesetzt werden.

Soll die Einstellung dauerhaft auch nach dem Ausschalten erhalten 
bleiben, muss die Konfiguration durch den Befehl \texttt{XW} gesichert werden.

Beispiele:

\begin{verbatim}
@XD=0
OPEN 1,8,15,"XD=0" : CLOSE 1
OPEN 1,8,15 : PRINT#1,"XD=";CHR$(0) : CLOSE 1
\end{verbatim}

\subsection{XI -- Init -- Einschaltkonfiguration wieder herstellen}
\index{XI}
\index{Einschaltkonfiguration}
\index{Konfiguration!wiederherstellen}
\index{Konfiguration!des Servers|see{Server}}

\texttt{XI}

Wenn der Adapter über einen nicht flüchtigen Speicher für die
Konfiguration besitzt, wird die dort gespeicherte Konfiguration beim
Einschalten des Adapters automatisch wieder hergestellt. 
Ist das nicht möglich, werden Standard-Werte eingesetzt.
Mit \texttt{XI} kann diese Einschaltkonfiguration jederzeit wieder
hergestellt werden.

Standardkonfiguration, falls keine gespeicherte Konfiguration verfügbar:

\begin{tabular}[c]{l l l}
\toprule 
Einstellung & Befehl & Wert\\
\midrule
Geräteadresse		& XU= & 8 \\
Letztes Laufwerk	& XD= & 0 \\
\bottomrule
\end{tabular}

\clearpage %MANFF
\subsection{XW -- Write configuration -- Konfiguration schreiben}

\texttt{XW}
\index{XW}
\index{Konfiguration!speichern}

Speichert die aktuelle Konfiguration in einem nicht flüchtigem Speicher.
Beim nächsten Einschalten des Adapters werden die gespeicherten Werte
automatisch wieder hergestellt. Eine Übersicht der gespeicherten
Einstellungen ist in der Be\-schrei\-bung des \texttt{XI}-Befehls enthalten.

\subsection{XRESET - Adapter zurücksetzen}
\index{XRESET}
\index{Reset}
\index{Bootloader}

\texttt{XRESET}

Dieser Befehl ist die Software-Entsprechung zum Druck auf die Reset-Taste
des Adapters. Alle nicht gespeicherten Daten und Einstellungen gehen dabei
verloren.

Die Daten des nichtflüchtigen Speichers bleiben erhalten. Die darin
gespeicherte Konfiguration kann sich jedoch von der zuletzt verwendeten
unterscheiden, wenn nicht zuvor ein \texttt{XW}-Befehl abgesetzt
wurde. 

Wenn noch Dateien zum Schreiben geöffnet sind, droht Datenverlust. 
Bei Dateisystemen, die auf dem Adapter gespeichert werden, wie etwa SD-Karten,
droht die Korrumpierung des gesamten Dateisystems, wenn zuvor
Schreibzugriffe erfolgt sind. Schließen Sie daher immer erst alle
geöffneten Dateien, bevor Sie diesen Befehl verwenden!

Wenn der Adapter über einen Bootloader zur Aktualisierung der Firmware
verfügt, wird dieser vor dem Starten der Firmware aktiv. \texttt{XRESET}
kann daher dazu verwendet werden, den Aktualisierungsvorgang anzustoßen.

\clearpage
\section{Konfiguration des Servers}
\label{serverconf}
Der Server kann durch folgende Kommandozeilenparameter konfiguriert 
werden:

\texttt{fsser} \Big[ \textit{-<Option> \ldots \ \Big] }
\textit{<Laufzeit-Verzeichnis>}

\subsection{Laufzeit-Verzeichnis}
\index{Laufzeit-Verzeichnis}
\index{Pfade}
Beim Starten wechselt der Server in das angegebene Unterverzeichnis
\textit{<Laufzeit-Verzeichnis>}. Alle in den Optionen angegebenen
relative Pfadangaben beziehen sich dann auf dieses Laufzeitverzeichnis.

Beispiel: Ihre CBM-Software liegt in einem Ordner mit dem Pfad

\texttt{/home/benutzer-name/Software/Commodore/CBM/}

Sie möchten nun die Laufwerke 0 bis 4 belegen.

Langform ohne relative Pfadangaben:

\begin{verbatim}
fsser 
-A0=fs:/home/benutzer-name/Software/Commodore/CBM/Visicalc
-A1=fs:/home/benutzer-name/Software/Commodore/CBM/Wordpro
-A2=fs:/home/benutzer-name/Software/Commodore/CBM/Assembler
-A3=fs:/home/benutzer-name/Software/Commodore/CBM/Spiele
-A4=fs:/home/benutzer-name/Software/Commodore/CBM/Spiele/Schach
.
\end{verbatim}

Kurzform mit Laufzeit-Verzeichnis und relativen Pfaden:

\begin{verbatim}
fsser 
-A0=fs:Visicalc
-A1=fs:Wordpro
-A2=fs:Assembler
-A3=fs:Spiele
-A4=fs:Spiele/Schach
/home/benutzer-name/Software/Commodore/CBM
\end{verbatim}

\subsection{Gerätedatei des Adapters}
\index{Ger\"atedatei}
\index{/dev/ttyUSB0}
\index{/dev/ttyAMA0}
\index{COM-Port}
\index{auto (Gerätedatei)}
\mbox{\Big[ \texttt{-d} \Big( \texttt{auto} \Big| 
\textit{<device>} \Big) \Big] }

Die Device-Option \texttt{-d} benennt die durch den
Adapter verwendete Gerätedatei, wie beispielsweise
\texttt{/dev/ttyUSB0}.

Wird diese Option nicht angegeben, wird 
\texttt{-d auto} verwendet.

Die automatische Erkennung \texttt{-d auto}
prüft, ob eine Reihe typischer Namen verfügbar sind, und wenn
genau einer dieser Möglichkeiten verfügbar ist, wird
diese verwendet. Findet der Server keine oder mehrere
Möglichkeiten wird er nicht starten und folgende Fehlermeldung
ausgeben:

\begin{verbatim}
Unable to decide which serial device it is.
Please pick one.
\end{verbatim}

Die Benennung unterscheidet sich je nach Betriebssystem.
Für den vom XS-1541 und petSD verwendeten Serial-over-USB Chip
FT232RL \index{FT232} weist Linux
fortlaufende Nummern beginnend bei \texttt{/dev/ttyUSB0} 
zu, während
\mbox{OS X} die im Chip gespei"-cherte Seriennummer verwendet. 
Diese Eigenart kann unter \mbox{OS X} dazu verwendet werden, die 
Seriennummer durch den Gerätenamen zu ersetzen.

\begin{tabular}[c]{l l l l}
\toprule 
Adapter & Betriebs- & Gerätedatei & Bemerkung\\
& system & & \\
\midrule
XS-1541 & Linux & /dev/ttyUSB0 & \\
XS-1541 & OS X & /dev/cu.usbserial-A10044Co & Standard-Seriennummer \\
petSD & Linux & /dev/ttyUSB0 & \\
petSD & OS X & /dev/cu.usbserial-petSD & Seriennummer bearbeitet \\
3.1541 & Linux & /dev/ttyAMA0 \\
\bottomrule
\end{tabular}

\subsection{Ausführlichkeit der Meldungen}
\index{verbose}
\index{Debug-Meldungen}
\mbox{\Big[ \texttt{-v} \Big]}

Wenn die Option \texttt{-v} \glqq{}verbose\grqq{} gegeben wird,
wird der Server eine Reihe weiterer ausführlicherer
Ausgaben machen, die bei der Fehlersuche behilflich sein können.

\subsection{Laufwerken Inhalte zuweisen}
\mbox{ \Big[ \texttt{-A}\textit{<Laufwerk-Nummer>} 
\texttt{:} \textit{<Provider>} \texttt{=}
\textit{<Parameter>} \Big] \Big[ \ldots \Big] }

Mit der \glqq{}assign\grqq{}-Option \texttt{-A} werden Laufwerke Inhalte
eines Providers zugewiesen. 

Einfach formuliert wird damit angegeben,
welche \glqq{}virtuelle Diskette\grqq{} in welches Laufwerk
eingeschoben wird. 

Provider können sowohl durch die XD-2031 Server Software,
die auf dem PC läuft, als auch durch den Adapter zur Verfügung
gestellt werden.

Neben dem Option \texttt{-A} beim
Serverstart können Zuweisungen auch vom CBM aus über das
Senden eines A-Befehls an den Kommandokanal \#15 erfolgen.  
Daher wird \textit{ASSIGN} in einem eigenem
Kapitel \ref{assign} \glqq{}\nameref{assign}\grqq{} 
auf Seite \pageref{assign} beschrieben.

\begin{comment}
TODO: is this feature broken?
Bei Providern, die der Server anbietet, ist zu beachten, 
dass aus Sicherheitsgründen zwar auf 
Unterverzeichnisse von \textit{<Parameter>} zugegriffen
werden kann, nicht jedoch auf Verzeichnisse in höherer Ebene.
\end{comment}

\subsection{Laufzeit-Konfiguration der Busse}
\index{Ger\"ateadresse!\"andern}
\mbox{\Big[ \texttt{-X} \textit{<Bus>} \texttt{:}
\textit{<Befehl>} \Big] \Big[ \ldots \Big] }

Mit Hilfe der X-Option können dem Adapter eine Reihe
X-Befehle übermittelt werden, die ab Seite \pageref{xcommands}
beschrieben werden.

Um beispielsweise die Geräteadresse für den IEC-Bus auf
9 und für den \mbox{IEEE-488} Bus auf 10 zu setzen:

\begin{verbatim}
-Xiec:U=9 -Xieee:U=10
\end{verbatim}

\clearpage
\section{Laufwerken Inhalte zuweisen}
\label{assign}
\index{ASSIGN}

\subsection{Allgemeines über Provider}

\mbox{ \Big[ \texttt{A}\textit{<Laufwerk-Nummer>} 
\texttt{:} \textit{<Provider>} \texttt{=}
\textit{<Parameter>} \Big] \Big[ \ldots \Big] }

Mit der \glqq{}assign\grqq{}-Option \texttt{-A} werden Laufwerke Inhalte
eines Providers zugewiesen. Einfach formuliert wird damit angegeben,
welche \glqq{}virtuelle Diskette\grqq{} in welches Laufwerk
eingeschoben wird.

\begin{tabular}[c]{l l}
\toprule 
Provider & Zugriff auf \\
\midrule
fs & Dateisystem des PC \\
di & Disk-Images im Dateisystem des PC  \\
ftp & FTP-Server (über PC, nur lesend) \\
http & Webserver (über PC, nur lesend) \\
\bottomrule
\end{tabular}

\subsection{fs: Zugriff auf das Dateisystem des PC}
\index{fs-Provider}
\mbox{ \texttt{A}\textit{<Laufwerk-Nummer>} 
\texttt{:} \textit{fs} \texttt{=}
\textit{<Pfad-im-PC-Dateisystem>} }

Der Filesystem-Provider bietet schreibenden und lesenden Zugriff auf
Verzeichnisse. Bei relativen Pfadangaben beziehen sich diese
auf das Laufzeit-Verzeichnis des Servers, bei absoluten Pfadangaben
spielt das Laufzeit-Verzeichnis natürlich keine Rolle.

In diesem Beispiel wird der absolute Pfad /usr/local/xd2031/tools
dem Laufwerk 0 zugewiesen:

\begin{verbatim}
fsser -A0:fs=/usr/local/xd2031/tools .
\end{verbatim}

Das Laufzeit-Verzeichnis des Servers wurde der Einfachheit halber
auf das aktuelle Verzeichnis '.' gesetzt.

Ein Beispiel mit relativen Pfaden:

\begin{verbatim}
fsser -A0:fs=office/text ~/Software/Commodore/CBM
\end{verbatim}

Pfadangaben sind relativ zum Laufzeit-Verzeichnis, 
so ergibt sich der Gesamt"-pfad  
\texttt{\~{}/Software/Commodore/CBM/office/text}.
Dieses Verzeichnis steht nun als Laufwerk 0 
für Schreib- und Lesevorgänge zur Verfügung.

\subsection{di: Zugriff auf Disk-Images im Dateisystem des PC}
\index{di-Provider}
\index{Disk Images}
\index{D64}
\index{D80}
\mbox{ \texttt{A}\textit{<Laufwerk-Nummer>} 
\texttt{:} \textit{di} \texttt{=}
\textit{<Disk-Image-im-PC-Dateisystem>} }

Mit dem Disk-Image-Provider können Disk-Images einem
Laufwerk zugewiesen werden. Ein Disk-Image enthält einen 1:1 Abzug
einer Diskette eines bestimmten Diskettenlaufwerks, inklusive
aller Daten des Dateisystems.

Die Erkennung des Typ des Disk-Images erfolgt unabhängig vom 
Dateinamen aufgrund der Dateigröße. 

Es gibt verschiedene Varianten, die sich durch die Anzahl
der hinterlegten Spuren (35 / 40) oder eine eventuell 
angehängte 
Fehlertabelle\footnote{Status-Code des Diskettencontrollers für 
jeden Block des Abbildes} 
unterscheiden können. Auch Kombinationen daraus sind möglich.

Varianten werden derzeit noch nicht unterstützt, lediglich die
Standard-Images.

Folgende Disk-Images werden unterstützt:

\begin{tabular}[c]{l l r}
\toprule 
Typ & Original-Laufwerke & Größe in Bytes \\
\midrule
D64 & 1541, 1571, 2031, 4031, 4040 & \numprint{174848} \\
D80 & 8050, 8250 & \numprint{533248} \\
D81 & 1581 & \numprint{819200} \\
D82 & 8250 & \numprint{1066496} \\
\bottomrule
\end{tabular}

Die Verwendung von Disk-Images bietet sich dort an,
wo die Erhaltung von Sonderzeichen im Dateinamen oder die
Anordnung innerhalb des Inhaltsverzeichnis wichtig ist.
\index{Sonderzeichen}
\index{PETSCII}
\index{Verzeichnis!Reihenfolge}
\index{Reihenfolge|see{Verzeichnis}}

Beispiel: die \glqq{}virtuelle Diskette\grqq{}
\texttt{visicalc.d82} wird in Laufwerk 0 eingelegt:

\begin{verbatim}
fsser -A0:di=office/calc/visicalc.d82 ~/Software/Commodore/CBM
\end{verbatim}

\clearpage
\section{Anhang}
\subsection{Verwendete Lizenzen}
\index{Lizenzen}
\subsubsection{Benutzerhandbuch}
Dieses Handbuch ist unter der
\textit{Creative Commons CC BY-SA 3.0} lizenziert.

\copyright{} 2013 Nils Eilers

Eine Beschreibung und den rechtsverbindlichen Lizenzvertrag
können Sie unter \url{http://creativecommons.org/licenses/by-sa/3.0/de/}
einsehen.

\includegraphics[scale=1]{by-sa.pdf}

\subsubsection{XD-2031}
Das Programm \glqq{}XD-2031\grqq{} ist unter der
\textit{GNU General Public License, Version 2} lizenziert

\copyright{} 2012-2013 André Fachat, Ingo Korb, Thomas Winkler,
Nils Eilers und andere.

Der rechtsverbindliche Lizenzvertrag ist im folgenden angehängt.
Bitte beachten Sie, dass der PC Server auch unter späteren
Versionen der GPL verfügbar ist, die ATmega Firmware dagegen
\textbf{nur unter GPL Version 2}. Die tatsächlich anwendbare
Lizenz ist im Quelltext der jeweiligen Datei enthalten.

\clearpage
\begin{center}

\mbox{\bf\huge GNU GENERAL PUBLIC LICENSE}

{\bf\large Version 2, June 1991}
\bigskip

{\parindent 0in

Copyright \copyright\ 1989, 1991 Free Software Foundation, Inc.

\bigskip

51 Franklin Street, Fifth Floor, Boston, MA  02110-1301, USA

\bigskip

Everyone is permitted to copy and distribute verbatim copies
of this license document, but changing it is not allowed.
}

{\bf\large Preamble}
\end{center}

The licenses for most software are designed to take away your freedom to
share and change it.  By contrast, the GNU General Public License is
intended to guarantee your freedom to share and change free software---to
make sure the software is free for all its users.  This General Public
License applies to most of the Free Software Foundation's software and to
any other program whose authors commit to using it.  (Some other Free
Software Foundation software is covered by the GNU Library General Public
License instead.)  You can apply it to your programs, too.

When we speak of free software, we are referring to freedom, not price.
Our General Public Licenses are designed to make sure that you have the
freedom to distribute copies of free software (and charge for this service
if you wish), that you receive source code or can get it if you want it,
that you can change the software or use pieces of it in new free programs;
and that you know you can do these things.

To protect your rights, we need to make restrictions that forbid anyone to
deny you these rights or to ask you to surrender the rights.  These
restrictions translate to certain responsibilities for you if you
distribute copies of the software, or if you modify it.

For example, if you distribute copies of such a program, whether gratis or
for a fee, you must give the recipients all the rights that you have.  You
must make sure that they, too, receive or can get the source code.  And
you must show them these terms so they know their rights.

We protect your rights with two steps: (1) copyright the software, and (2)
offer you this license which gives you legal permission to copy,
distribute and/or modify the software.

Also, for each author's protection and ours, we want to make certain that
everyone understands that there is no warranty for this free software.  If
the software is modified by someone else and passed on, we want its
recipients to know that what they have is not the original, so that any
problems introduced by others will not reflect on the original authors'
reputations.

Finally, any free program is threatened constantly by software patents.
We wish to avoid the danger that redistributors of a free program will
individually obtain patent licenses, in effect making the program
proprietary.  To prevent this, we have made it clear that any patent must
be licensed for everyone's free use or not licensed at all.

The precise terms and conditions for copying, distribution and
modification follow.

\begin{center}
{\Large \sc Terms and Conditions For Copying, Distribution and
  Modification}
\end{center}


%\renewcommand{\theenumi}{\alpha{enumi}}
\begin{enumerate}

\addtocounter{enumi}{-1}

\item 

This License applies to any program or other work which contains a notice
placed by the copyright holder saying it may be distributed under the
terms of this General Public License.  The ``Program'', below, refers to
any such program or work, and a ``work based on the Program'' means either
the Program or any derivative work under copyright law: that is to say, a
work containing the Program or a portion of it, either verbatim or with
modifications and/or translated into another language.  (Hereinafter,
translation is included without limitation in the term ``modification''.)
Each licensee is addressed as ``you''.

Activities other than copying, distribution and modification are not
covered by this License; they are outside its scope.  The act of
running the Program is not restricted, and the output from the Program
is covered only if its contents constitute a work based on the
Program (independent of having been made by running the Program).
Whether that is true depends on what the Program does.

\item You may copy and distribute verbatim copies of the Program's source
  code as you receive it, in any medium, provided that you conspicuously
  and appropriately publish on each copy an appropriate copyright notice
  and disclaimer of warranty; keep intact all the notices that refer to
  this License and to the absence of any warranty; and give any other
  recipients of the Program a copy of this License along with the Program.

You may charge a fee for the physical act of transferring a copy, and you
may at your option offer warranty protection in exchange for a fee.

\item

You may modify your copy or copies of the Program or any portion
of it, thus forming a work based on the Program, and copy and
distribute such modifications or work under the terms of Section 1
above, provided that you also meet all of these conditions:

\begin{enumerate}

\item 

You must cause the modified files to carry prominent notices stating that
you changed the files and the date of any change.

\item

You must cause any work that you distribute or publish, that in
whole or in part contains or is derived from the Program or any
part thereof, to be licensed as a whole at no charge to all third
parties under the terms of this License.

\item
If the modified program normally reads commands interactively
when run, you must cause it, when started running for such
interactive use in the most ordinary way, to print or display an
announcement including an appropriate copyright notice and a
notice that there is no warranty (or else, saying that you provide
a warranty) and that users may redistribute the program under
these conditions, and telling the user how to view a copy of this
License.  (Exception: if the Program itself is interactive but
does not normally print such an announcement, your work based on
the Program is not required to print an announcement.)

\end{enumerate}


These requirements apply to the modified work as a whole.  If
identifiable sections of that work are not derived from the Program,
and can be reasonably considered independent and separate works in
themselves, then this License, and its terms, do not apply to those
sections when you distribute them as separate works.  But when you
distribute the same sections as part of a whole which is a work based
on the Program, the distribution of the whole must be on the terms of
this License, whose permissions for other licensees extend to the
entire whole, and thus to each and every part regardless of who wrote it.

Thus, it is not the intent of this section to claim rights or contest
your rights to work written entirely by you; rather, the intent is to
exercise the right to control the distribution of derivative or
collective works based on the Program.

In addition, mere aggregation of another work not based on the Program
with the Program (or with a work based on the Program) on a volume of
a storage or distribution medium does not bring the other work under
the scope of this License.

\item
You may copy and distribute the Program (or a work based on it,
under Section 2) in object code or executable form under the terms of
Sections 1 and 2 above provided that you also do one of the following:

\begin{enumerate}

\item

Accompany it with the complete corresponding machine-readable
source code, which must be distributed under the terms of Sections
1 and 2 above on a medium customarily used for software interchange; or,

\item

Accompany it with a written offer, valid for at least three
years, to give any third party, for a charge no more than your
cost of physically performing source distribution, a complete
machine-readable copy of the corresponding source code, to be
distributed under the terms of Sections 1 and 2 above on a medium
customarily used for software interchange; or,

\item

Accompany it with the information you received as to the offer
to distribute corresponding source code.  (This alternative is
allowed only for noncommercial distribution and only if you
received the program in object code or executable form with such
an offer, in accord with Subsection b above.)

\end{enumerate}


The source code for a work means the preferred form of the work for
making modifications to it.  For an executable work, complete source
code means all the source code for all modules it contains, plus any
associated interface definition files, plus the scripts used to
control compilation and installation of the executable.  However, as a
special exception, the source code distributed need not include
anything that is normally distributed (in either source or binary
form) with the major components (compiler, kernel, and so on) of the
operating system on which the executable runs, unless that component
itself accompanies the executable.

If distribution of executable or object code is made by offering
access to copy from a designated place, then offering equivalent
access to copy the source code from the same place counts as
distribution of the source code, even though third parties are not
compelled to copy the source along with the object code.

\item
You may not copy, modify, sublicense, or distribute the Program
except as expressly provided under this License.  Any attempt
otherwise to copy, modify, sublicense or distribute the Program is
void, and will automatically terminate your rights under this License.
However, parties who have received copies, or rights, from you under
this License will not have their licenses terminated so long as such
parties remain in full compliance.

\item
You are not required to accept this License, since you have not
signed it.  However, nothing else grants you permission to modify or
distribute the Program or its derivative works.  These actions are
prohibited by law if you do not accept this License.  Therefore, by
modifying or distributing the Program (or any work based on the
Program), you indicate your acceptance of this License to do so, and
all its terms and conditions for copying, distributing or modifying
the Program or works based on it.

\item
Each time you redistribute the Program (or any work based on the
Program), the recipient automatically receives a license from the
original licensor to copy, distribute or modify the Program subject to
these terms and conditions.  You may not impose any further
restrictions on the recipients' exercise of the rights granted herein.
You are not responsible for enforcing compliance by third parties to
this License.

\item
If, as a consequence of a court judgment or allegation of patent
infringement or for any other reason (not limited to patent issues),
conditions are imposed on you (whether by court order, agreement or
otherwise) that contradict the conditions of this License, they do not
excuse you from the conditions of this License.  If you cannot
distribute so as to satisfy simultaneously your obligations under this
License and any other pertinent obligations, then as a consequence you
may not distribute the Program at all.  For example, if a patent
license would not permit royalty-free redistribution of the Program by
all those who receive copies directly or indirectly through you, then
the only way you could satisfy both it and this License would be to
refrain entirely from distribution of the Program.

If any portion of this section is held invalid or unenforceable under
any particular circumstance, the balance of the section is intended to
apply and the section as a whole is intended to apply in other
circumstances.

It is not the purpose of this section to induce you to infringe any
patents or other property right claims or to contest validity of any
such claims; this section has the sole purpose of protecting the
integrity of the free software distribution system, which is
implemented by public license practices.  Many people have made
generous contributions to the wide range of software distributed
through that system in reliance on consistent application of that
system; it is up to the author/donor to decide if he or she is willing
to distribute software through any other system and a licensee cannot
impose that choice.

This section is intended to make thoroughly clear what is believed to
be a consequence of the rest of this License.

\item
If the distribution and/or use of the Program is restricted in
certain countries either by patents or by copyrighted interfaces, the
original copyright holder who places the Program under this License
may add an explicit geographical distribution limitation excluding
those countries, so that distribution is permitted only in or among
countries not thus excluded.  In such case, this License incorporates
the limitation as if written in the body of this License.

\item
The Free Software Foundation may publish revised and/or new versions
of the General Public License from time to time.  Such new versions will
be similar in spirit to the present version, but may differ in detail to
address new problems or concerns.

Each version is given a distinguishing version number.  If the Program
specifies a version number of this License which applies to it and ``any
later version'', you have the option of following the terms and conditions
either of that version or of any later version published by the Free
Software Foundation.  If the Program does not specify a version number of
this License, you may choose any version ever published by the Free Software
Foundation.

\item
If you wish to incorporate parts of the Program into other free
programs whose distribution conditions are different, write to the author
to ask for permission.  For software which is copyrighted by the Free
Software Foundation, write to the Free Software Foundation; we sometimes
make exceptions for this.  Our decision will be guided by the two goals
of preserving the free status of all derivatives of our free software and
of promoting the sharing and reuse of software generally.

\begin{center}
{\Large\sc
No Warranty
}
\end{center}

\item
{\sc Because the program is licensed free of charge, there is no warranty
for the program, to the extent permitted by applicable law.  Except when
otherwise stated in writing the copyright holders and/or other parties
provide the program ``as is'' without warranty of any kind, either expressed
or implied, including, but not limited to, the implied warranties of
merchantability and fitness for a particular purpose.  The entire risk as
to the quality and performance of the program is with you.  Should the
program prove defective, you assume the cost of all necessary servicing,
repair or correction.}

\item
{\sc In no event unless required by applicable law or agreed to in writing
will any copyright holder, or any other party who may modify and/or
redistribute the program as permitted above, be liable to you for damages,
including any general, special, incidental or consequential damages arising
out of the use or inability to use the program (including but not limited
to loss of data or data being rendered inaccurate or losses sustained by
you or third parties or a failure of the program to operate with any other
programs), even if such holder or other party has been advised of the
possibility of such damages.}

\end{enumerate}


\begin{center}
{\Large\sc End of Terms and Conditions}
\end{center}


\pagebreak[2]

\section*{Appendix: How to Apply These Terms to Your New Programs}

If you develop a new program, and you want it to be of the greatest
possible use to the public, the best way to achieve this is to make it
free software which everyone can redistribute and change under these
terms.

  To do so, attach the following notices to the program.  It is safest to
  attach them to the start of each source file to most effectively convey
  the exclusion of warranty; and each file should have at least the
  ``copyright'' line and a pointer to where the full notice is found.

\begin{quote}
one line to give the program's name and a brief idea of what it does. \\
Copyright (C) yyyy  name of author \\

This program is free software; you can redistribute it and/or modify
it under the terms of the GNU General Public License as published by
the Free Software Foundation; either version 2 of the License, or
(at your option) any later version.

This program is distributed in the hope that it will be useful,
but WITHOUT ANY WARRANTY; without even the implied warranty of
MERCHANTABILITY or FITNESS FOR A PARTICULAR PURPOSE.  See the
GNU General Public License for more details.

You should have received a copy of the GNU General Public License
along with this program; if not, write to the Free Software
Foundation, Inc., 51 Franklin Street, Fifth Floor, Boston, MA  02110-1301, USA.
\end{quote}

Also add information on how to contact you by electronic and paper mail.

If the program is interactive, make it output a short notice like this
when it starts in an interactive mode:

\begin{quote}
Gnomovision version 69, Copyright (C) yyyy  name of author \\
Gnomovision comes with ABSOLUTELY NO WARRANTY; for details type `show w'. \\
This is free software, and you are welcome to redistribute it
under certain conditions; type `show c' for details.
\end{quote}


The hypothetical commands {\tt show w} and {\tt show c} should show the
appropriate parts of the General Public License.  Of course, the commands
you use may be called something other than {\tt show w} and {\tt show c};
they could even be mouse-clicks or menu items---whatever suits your
program.

You should also get your employer (if you work as a programmer) or your
school, if any, to sign a ``copyright disclaimer'' for the program, if
necessary.  Here is a sample; alter the names:

\begin{quote}
Yoyodyne, Inc., hereby disclaims all copyright interest in the program \\
`Gnomovision' (which makes passes at compilers) written by James Hacker. \\

signature of Ty Coon, 1 April 1989 \\
Ty Coon, President of Vice
\end{quote}


This General Public License does not permit incorporating your program
into proprietary programs.  If your program is a subroutine library, you
may consider it more useful to permit linking proprietary applications
with the library.  If this is what you want to do, use the GNU Library
General Public License instead of this License.

\clearpage


\renewcommand{\indexname}{Stichwortverzeichnis}
\printindex

\end{document}




